\chapter{Evaluasi Model yang Jujur}

Sebuah laporan pendek sampai di meja seorang analis keamanan. Isinya satu tabel: detektor
intrusi baru, diuji pada data lalu lintas jaringan yang diambil dari jaringan kampus
selama sebulan, mencapai akurasi $99$ persen. Angka itu dibingkai dengan huruf tebal dan
satu kalimat penutup: ``siap dipakai''.

Data itu berisi satu serangan untuk setiap seribu koneksi. Pada kumpulan yang sama,
program tiga baris yang selalu menjawab ``bersih'' memperoleh akurasi $99{,}9$ persen: ia
tidak membaca satu fitur pun dan tidak pernah salah menuduh. Detektor yang dilatih berjam-
jam itu kalah, dan laporannya tidak memuat satu angka pun yang membocorkan kekalahan itu.

Yang keliru adalah ringkasannya: satu bilangan mewakili empat bilangan yang perannya
berbeda, pada data yang komposisinya tidak disebut. Bab 26 menjelaskan mengapa galat uji
berbeda dari galat latih; bab ini menutup tiga celah: metrik apa yang jujur tentang apa
yang mereka abaikan, seberapa besar angkanya bergoyang, dan apakah data uji benar-benar
belum dilihat oleh prosedur yang memilih model.

\begin{untukanak}
Bayangkan sebuah detektor logam di gerbang sekolah. Setiap pagi ia berbunyi. Kemarin ia
berbunyi seribu kali, dan petugas memeriksa seribu tas. Ternyata hanya enam puluh tas
yang benar-benar membawa barang terlarang. Sembilan ratus empat puluh siswa lain harus
membuka tasnya tanpa alasan.

Sekarang datang petugas dinas yang bertanya: berapa persen jawaban alat itu yang benar?
Petugas gerbang menjawab dengan bangga: sembilan puluh sembilan persen. Ia tidak
berbohong; dari seratus ribu kantong, alat itu hanya salah pada seribu. Tetapi jawaban itu
tidak menjawab pertanyaan yang membuat siswa kesal: mereka kesal bukan karena alat itu
sering salah, melainkan karena hampir setiap kali berbunyi, ia salah.

Dua pertanyaan itu terdengar mirip dan jawabannya berbeda jauh: yang satu bertanya tentang
seluruh kantong, yang lain tentang kantong yang dibunyikan. Di bab ini kita akan tahu
mengapa satu alat bisa unggul pada pertanyaan pertama dan hancur pada pertanyaan kedua,
mengapa angka seratus persen kadang bukan kabar baik, dan mengapa laporan yang jujur selalu
memuat lebih dari satu angka.
\end{untukanak}

\section{Empat bilangan yang tidak simetris}
\tingkat{dasar}

Pengklasifikasi biner, secanggih apa pun isinya, meninggalkan jejak yang sama: empat
bilangan. Yang membedakan satu laporan evaluasi dari laporan lain adalah apakah keempatnya
dilaporkan. Kalimat sehari-harinya: setiap keputusan punya dua cara untuk benar dan dua
cara untuk salah, dan menutup salah satunya sama dengan bersembunyi.

\kamus{matriks konfusi}{confusion matrix}{buku kas empat kolom yang mencatat setiap
keputusan benar dan salah beserta arah kesalahannya}

\begin{definisi}[Matriks konfusi]
\index{matriks konfusi}
Misalkan setiap contoh memiliki label sebenarnya $y\in\{0,1\}$ dan prediksi
$\hat y\in\{0,1\}$, dengan $1$ berarti positif (serangan, penyakit, pesan berbahaya).
Empat bilangan
\begin{align}
\mathrm{TP}&=\#\{y=1,\ \hat y=1\}, & \mathrm{FP}&=\#\{y=0,\ \hat y=1\},
\nonumber\\
\mathrm{FN}&=\#\{y=1,\ \hat y=0\}, & \mathrm{TN}&=\#\{y=0,\ \hat y=0\}
\nonumber
\end{align}
beserta susunannya disebut \emph{matriks konfusi}. Jumlah seluruh contoh adalah
$n=\mathrm{TP}+\mathrm{FP}+\mathrm{FN}+\mathrm{TN}$.
\end{definisi}

Dari empat bilangan itu diturunkan metrik yang menjawab pertanyaan berbeda. Presisi
bertanya: bila model mengatakan positif, berapa peluang ia benar. Recall bertanya: dari
seluruh kasus positif, berapa bagian yang tertangkap.

\begin{definisi}[Presisi, recall, spesifisitas]
\index{presisi}
\index{recall}
Dengan $\mathrm{TP}+\mathrm{FP}>0$ dan $\mathrm{TP}+\mathrm{FN}>0$,
\[
\operatorname{prec}=\frac{\mathrm{TP}}{\mathrm{TP}+\mathrm{FP}},\qquad
\operatorname{rec}=\frac{\mathrm{TP}}{\mathrm{TP}+\mathrm{FN}},\qquad
\operatorname{spec}=\frac{\mathrm{TN}}{\mathrm{TN}+\mathrm{FP}} .
\]
Akurasi adalah $(\mathrm{TP}+\mathrm{TN})/n$, dan $F_1$ adalah
$2\,\mathrm{TP}/(2\,\mathrm{TP}+\mathrm{FP}+\mathrm{FN})$.
\end{definisi}

Mari kita pakai satu contoh kerja: sebuah detektor diuji pada $1000$ koneksi berlabel,
$100$ di antaranya serangan. Ia menandai $62$ serangan dengan benar, melewatkan
$38$, dan menuduh $17$ koneksi bersih sebagai serangan.

\begin{center}
\small
\begin{tabular}{l|rr|r}
\hline
 & prediksi positif & prediksi negatif & total \\
\hline
sebenarnya positif & $\mathrm{TP}=62$ & $\mathrm{FN}=38$ & 100 \\
sebenarnya negatif & $\mathrm{FP}=17$ & $\mathrm{TN}=883$ & 900 \\
\hline
total & 79 & 921 & 1000 \\
\hline
\end{tabular}
\end{center}

Presisinya $62/79=0{,}784810$: sekitar empat dari lima alarm menunjuk serangan.
Recall-nya $62/100=0{,}620000$: hampir empat dari sepuluh serangan tidak terlihat.
Spesifisitasnya $883/900=0{,}981111$. Akurasinya $945/1000=0{,}945000$, dan
$F_1=124/179=0{,}692737$. Kelima angka itu benar tentang model yang sama, dan tidak satu
pun darinya boleh dipakai sendiri. Pemeriksaan dengan \texttt{sklearn.metrics} memberi
kelima angka yang sama sampai enam angka desimal.

Gambar~\ref{fig:konfusi-jujur} menempatkan keempat bilangan itu pada tempatnya, lengkap
dengan jumlah baris dan kolomnya; setiap angka di sisa bab harus dapat dilacak ke salah
satu sel gambar itu.

\begin{figure}[htbp]
\centering
\begin{tikzpicture}[every node/.style={font=\footnotesize}]
  % dua sel benar (hijau) dan dua sel salah (merah)
  \fill[layerjade!14] (4.40,2.00) rectangle (6.70,2.90);
  \fill[layerjade!14] (6.70,1.10) rectangle (9.00,2.00);
  \fill[accent!8]     (6.70,2.00) rectangle (9.00,2.90);
  \fill[accent!8]     (4.40,1.10) rectangle (6.70,2.00);
  \draw[ink,thick] (4.40,0.55) rectangle (10.30,2.90);
  \draw[ink] (6.70,0.55) -- (6.70,2.90);
  \draw[ink] (9.00,0.55) -- (9.00,2.90);
  \draw[ink] (4.40,2.00) -- (10.30,2.00);
  \draw[ink] (4.40,1.10) -- (10.30,1.10);
  \draw[ink!35] (4.40,0.55) -- (10.30,0.55);
  % isi sel
  \node at (5.55,2.45) {$\mathrm{TP}=62$};
  \node at (7.85,2.45) {$\mathrm{FN}=38$};
  \node at (5.55,1.55) {$\mathrm{FP}=17$};
  \node at (7.85,1.55) {$\mathrm{TN}=883$};
  % judul kolom
  \node at (5.55,3.14) {prediksi positif};
  \node at (7.85,3.14) {prediksi negatif};
  \node at (9.65,3.14) {total};
  % judul baris
  \node[anchor=east] at (4.24,2.45) {sebenarnya positif $(100)$};
  \node[anchor=east] at (4.24,1.55) {sebenarnya negatif $(900)$};
  \node[anchor=east] at (4.24,0.82) {total};
  % jumlah
  \node at (5.55,0.82) {$79$};
  \node at (7.85,0.82) {$921$};
  \node at (9.65,2.45) {$100$};
  \node at (9.65,1.55) {$900$};
  \node at (9.65,0.82) {$1000$};
  \node[anchor=north,font=\scriptsize,align=center,text width=11.4cm] at (7.00,0.30)
    {akurasi $=(62+883)/1000=0{,}945$;\quad presisi $=62/79$;\quad recall $=62/100$};
\end{tikzpicture}
\caption{Matriks konfusi contoh kerja pada $1000$ koneksi: $62$ serangan tertangkap, $38$
lolos, $17$ alarm palsu, dan $883$ koneksi bersih yang ditolak dengan benar. Baris
menjumlah menjadi $100$ positif dan $900$ negatif, kolom menjadi $79$ prediksi positif dan
$921$ prediksi negatif, dan seluruhnya $1000$. Dua sel hijau adalah dua cara untuk benar,
dua sel merah dua cara untuk salah. Akurasi $0{,}945$ menjumlahkan kedua sel hijau dan
karena itu tidak dapat membedakan $62$ yang tertangkap dari $38$ yang lolos.}
\label{fig:konfusi-jujur}
\end{figure}

Peran keempat bilangan tidak simetris. Presisi menjawab pertanyaan operator: bila
alarm berbunyi, berapa peluang memang ada serangan. Recall menjawab pertanyaan
perancang kebijakan: dari seluruh serangan, berapa bagian yang tertangkap. Keduanya
bergantung pada ambang, sehingga metrik tanpa ambang bukan pernyataan yang lengkap:
model yang sama memberi presisi $0{,}25$ dengan recall $0{,}45$, atau presisi
$0{,}013$ dengan recall $0{,}86$.

$F_1$ mencoba meringkas presisi dan recall menjadi satu bilangan, dan justru di situ ia
paling mudah disalahgunakan.

\begin{proposisi}
\index{f1}
$F_1$ adalah rata-rata harmonik presisi $P$ dan recall $R$, yaitu $F_1=2PR/(P+R)$,
sehingga $F_1\le(P+R)/2$ dengan kesamaan berlaku tepat ketika $P=R$.
\end{proposisi}

\begin{proof}[Bukti]
Tulis $a=\mathrm{TP}$, $b=\mathrm{FP}$, $c=\mathrm{FN}$. Maka $P=a/(a+b)$,
$R=a/(a+c)$, dan
\[
\frac{2PR}{P+R}
=\frac{2a^2\big/\bigl((a+b)(a+c)\bigr)}{a(2a+b+c)\big/\bigl((a+b)(a+c)\bigr)}
=\frac{2a}{2a+b+c}=F_1 .
\]
Ketaksamaan harmonik--aritmetik $4PR\le(P+R)^2$ ekuivalen dengan $(P-R)^2\ge0$, dan
kesamaan hanya terjadi bila $P=R$.
\end{proof}

Sifat harmonik itu menarik nilainya ke bawah oleh komponen terkecil: presisi $0{,}9$
dengan recall $0{,}2$ memberi $F_1\approx0{,}327$, bukan $0{,}55$; pada contoh kita
$F_1=0{,}6927$, sedikit di bawah titik tengah $0{,}7024$.

Tiga cacat ringkasan perlu dinamai satu per satu, sebab obatnya berbeda. Presisi
buta terhadap serangan yang lolos: model yang menandai semua koneksi memperoleh
presisi $0{,}1$, yaitu kelazimannya sendiri, tetapi recall $1$. $F_1$ mengabaikan
$\mathrm{TN}$ sepenuhnya: dua model dengan presisi dan recall sama tetapi alarm
palsu berbeda sepuluh kali lipat melaporkan $F_1$ identik, padahal biaya
operasionalnya jauh berbeda. $F_1$ juga membebani kedua komponen setara, padahal
satu serangan yang lolos biasanya jauh lebih mahal daripada satu alarm palsu.

\begin{tantangan}
Sebuah laporan hanya memuat tiga angka: ``akurasi $0{,}9720$ pada $5000$ koneksi uji;
presisi $0{,}000$; model siap dipakai''. Tentukan seluruh matriks konfusi yang konsisten
dengan laporan itu, jelaskan mengapa recall-nya nol dan bukan bilangan lain, lalu cari
kelaziman terbesar yang masih mungkin.

Petunjuk 1: akurasi hanya menjumlahkan dua sel diagonal, jadi
$5000(1-0{,}9720)=140$ kesalahan bebas dibagi antara dua sel yang tersisa.

Petunjuk 2: presisi $0{,}000$ berarti $\mathrm{TP}=0$, sehingga $\mathrm{FN}$ sama dengan
banyaknya contoh positif dan $\mathrm{FP}$ adalah sisanya. Tetapi presisi dengan
$\mathrm{TP}=\mathrm{FP}=0$ adalah $0/0$, bukan $0$; angka $0{,}000$ pada laporan itu
karena itu menuntut sekurang-kurangnya satu alarm palsu.

Kalau mentok, kembali ke bagian Akurasi yang menipu pada bab ini.
\end{tantangan}

\section{Akurasi yang menipu pada data tak seimbang}
\tingkat{dasar}

Akurasi paling sering dilaporkan dan paling mudah menipu: ia menghitung semua
kesalahan dengan harga sama pada kumpulan yang komposisinya tidak disebut. Kalimat
sehari-harinya: akurasi adalah nilai rapor yang dihitung seolah kedua kelompok
seimbang.

\index{akurasi seimbang}
\begin{definisi}[Akurasi seimbang]
Misalkan $\pi$ adalah kelaziman kelas positif. \emph{Akurasi seimbang} adalah
$(\operatorname{rec}+\operatorname{spec})/2$, yaitu rata-rata laju benar pada kedua kelas.
Besaran ini tidak berubah bila jumlah contoh positif dan negatif diubah tanpa mengubah
kedua laju tersebut.
\end{definisi}

Akurasi biasa tidak memiliki sifat itu, dan akibatnya dapat ditulis tepat.

\begin{proposisi}
Pengklasifikasi yang selalu menjawab negatif memperoleh akurasi $1-\pi$. Karena itu
sebuah model dengan akurasi di bawah $1-\pi$ kalah dari garis dasar yang tidak belajar
apa pun, dan laporan yang menyebut akurasi tanpa menyebut $\pi$ tidak dapat dinilai.
\end{proposisi}

\begin{proof}[Bukti]
Dari $n$ contoh, $\pi n$ berlabel positif dan seluruhnya salah, sedangkan $(1-\pi)n$
berlabel negatif dan seluruhnya benar. Akurasinya $(1-\pi)n/n=1-\pi$. Model lain dengan
akurasi lebih kecil karena itu melakukan lebih banyak kesalahan daripada garis dasar itu.
\end{proof}

\begin{bacakeras}
Dua suara, dibaca bergantian dengan teman. Ketuk meja setiap kali kata ``nol'' diucapkan.

``Seratus kartu. Berapa merah?'' ``Sepuluh.'' ``Berapa putih?'' ``Sembilan puluh.''
``Mesin menjawab apa?'' ``Putih. Semua putih.'' ``Berapa yang benar?'' ``Sembilan puluh.''
``Jadi akurasinya?'' ``Sembilan puluh persen.'' ``Dan yang merah?'' ``Nol. Tidak satu pun.''
``Angka tadi bohong?'' ``Tidak. Ia hanya menjawab satu pertanyaan.''
``Pertanyaan apa?'' ``Berapa banyak yang kutebak benar, dari semua yang ada.''
``Pertanyaan mana yang tidak dijawab?'' ``Dari yang benar-benar berbahaya, berapa yang kutemukan.''

Ulangi tiga baris tengah, makin cepat, dengan ketukan yang sama: yang dihitung semua, yang
hilang tak disebut, yang hilang tak disebut.

Suara kedua menutup pelan: kelas yang sedikit mudah tenggelam di dalam rata-rata kelas
yang banyak, dan angka rata-rata itu tidak pernah menyebut siapa yang tenggelam.

Baris yang harus diulang pembaca, sekali, dengan suara jelas: akurasi menjumlahkan yang
benar, bukan yang penting, jadi pada kelas tak seimbang sembilan puluh persen bisa berarti
nol.
\end{bacakeras}

Pada contoh kerja kita $\pi=0{,}1$, jadi garis dasarnya $0{,}9$ dan model dengan
$0{,}945$ memang lebih baik. Sekarang ubah kelazimannya saja, tanpa mengubah mutu model,
dan lihat apa yang terjadi pada angka yang sama.

\begin{contoh}[Satu serangan untuk setiap seribu koneksi]
Ambil $100\,000$ koneksi dengan $100$ serangan. Skor pada koneksi bersih menyebar
normal dengan rataan $0$ dan simpangan baku $1$; skor pada serangan dengan rataan
$2{,}6$ dan simpangan baku $1$. Ambang dipasang pada $2{,}3263$, yaitu
kuantil $0{,}99$ dari sebaran skor negatif, sehingga secara rancangan sekitar satu persen
koneksi bersih akan ditandai. Hasil satu lintasan dengan benih tetap:
\[
\mathrm{TP}=62,\quad \mathrm{FP}=1013,\quad \mathrm{FN}=38,\quad \mathrm{TN}=98\,887 .
\]
Akurasinya $(62+98\,887)/100\,000=0{,}989490$. Garis dasarnya
$99\,900/100\,000=0{,}999000$. Presisinya $62/1075=0{,}057674$, recall
$62/100=0{,}620000$, dan akurasi seimbangnya $0{,}804930$.
\end{contoh}

Angka $0{,}9895$ itu, bila dibulatkan, menjadi ``$99$ persen''. Ia bukan hanya tidak
informatif: ia lebih rendah daripada garis dasar sebesar $0{,}009510$, hampir satu poin
persentase penuh. Model itu melakukan sepuluh kali lebih banyak kesalahan daripada program
yang selalu menjawab ``bersih'', $1051$ melawan $100$, dan seluruhnya berupa alarm palsu,
$1013$ buah atau sekitar $16{,}3$ per serangan yang tertangkap.

Di sini pembulatan menambah kerusakan: $0{,}9895$ dan $0{,}9990$ sama-sama menjadi
``$99$ persen'', padahal keduanya terpisah oleh seluruh selisih yang menentukan.
Laporan yang baik karena itu menuliskan tiga atau empat angka desimal, menyebut $n$,
dan menyebut kelazimannya; tanpa ketiganya, pembaca tidak tahu apakah $99$ persen
itu prestasi atau kegagalan.

Apakah akurasi lalu tidak boleh dipakai? Ia sah bila biaya kedua kesalahan setara,
kelaziman dekat setengah, dan ambangnya tetap. Di luar itu ia harus ditemani presisi dan
recall, atau diganti akurasi seimbang: $0{,}8049$ lebih jujur daripada $0{,}9895$, sebab
kelas minoritas hanya tertangkap $62$ persen.

\begin{tangan}
Bahan: $100$ kartu, $90$ bermuka putih dan $10$ bermuka merah, spidol, dan satu lembar
tally dua kolom.

Tulis ``bersih'' pada sembilan puluh kartu putih dan ``serangan'' pada sepuluh kartu merah.
Kocok, lalu bagi menjadi sepuluh tumpukan. Aturan pertama: setiap tumpukan dijawab
``bersih'' tanpa membalik satu kartu pun. Catat jawaban yang benar, lalu balik semua kartu:
sembilan puluh benar, sepuluh merah lolos. Tulis di papan: akurasi
sembilan puluh persen, recall nol persen.

Sekarang aturan kedua. Beri titik di sudut enam kartu merah, dan tanpa memberi tahu pemain,
beri titik yang sama di sudut empat kartu putih. Jawab ``serangan'' hanya bila titik itu
terlihat. Hitung ulang: enam merah tertangkap, empat merah lolos, empat putih dituduh,
delapan puluh enam putih benar.

Yang harus terlihat: akurasi hanya bergerak dari sembilan puluh ke sembilan puluh dua
persen, sedangkan recall melompat dari nol ke enam puluh persen. Angka yang berubah paling
banyak justru angka yang tidak ada pada laporan pertama.
\end{tangan}

\section{Kurva ROC, AUC, dan prevalensi}
\tingkat{menengah}

Setiap ambang memberi satu titik pada dua laju yang berbeda; memilih ambang tanpa melihat
seluruh kemungkinannya adalah menebak. Kalimat sehari-harinya: kita menyapu ambang
dari longgar ke ketat, lalu menggambar jejak yang ditinggalkan model.

\kamus{kurva roc}{receiver operating characteristic}{jejak yang ditinggalkan sebuah model
ketika ambangnya disapu dari longgar ke ketat, digambar pada dua laju yang tidak
bergantung pada kelaziman}

\begin{definisi}[Kurva ROC dan AUC]
\index{kurva roc}
\index{auc}
Untuk ambang $t$, laju benar positif dan laju salah positif adalah
\[
\mathrm{TPR}(t)=P(S>t\mid y=1),\qquad \mathrm{FPR}(t)=P(S>t\mid y=0),
\]
dengan $S$ skor model. Kurva ROC adalah himpunan titik $(\mathrm{FPR}(t),\mathrm{TPR}(t))$
untuk seluruh $t$. Luas di bawah kurva itu disebut AUC, dan untuk data terhingga ia
dihitung sebagai luas trapesium di antara titik-titik yang teramati.
\end{definisi}

AUC punya arti yang jauh lebih sederhana daripada namanya, dan arti itu yang membuatnya
populer.

\begin{proposisi}
Untuk data tanpa skor kembar, AUC sama dengan peluang bahwa skor satu contoh positif
yang diambil acak melebihi skor satu contoh negatif yang diambil acak. Bila ada skor
kembar, tiap pasangan kembar dihitung setengah.
\end{proposisi}

\begin{proof}[Bukti]
Susun seluruh skor dari besar ke kecil dan sapu ambang ke bawah. Setiap kali ambang
melewati satu contoh positif, $\mathrm{TPR}$ naik sebesar $1/n_+$ pada ketinggian
$\mathrm{FPR}$ saat itu, yaitu bagian contoh negatif yang sudah terlewati. Setiap kali
ambang melewati contoh negatif, $\mathrm{FPR}$ naik pada $\mathrm{TPR}$ tetap dan tidak
menambah luas. Karena itu luas totalnya
\[
\mathrm{AUC}=\frac1{n_+}\sum_{i:\ y_i=1}\mathrm{FPR}_i
=\frac{1}{n_+n_-}\#\{(i,j):\ y_i=1,\ y_j=0,\ S_i>S_j\},
\]
yaitu peluang yang diklaim, dan suku terakhir itulah yang dihitung oleh statistik
Mann--Whitney. Skor kembar tidak menambah luas, sehingga tiap pasangan seperti itu
menyumbang setengah.
\end{proof}

AUC tidak melihat kelaziman: bila setiap contoh negatif digandakan sepuluh kali, urutan
skor tidak berubah dan AUC tetap. Cara menghitungnya di kertas diperlihatkan pada bagian
Contoh hitung.

Sekarang bagian yang membuat ROC berbahaya: sapu ambang simulasi satu banding seribu.

\begin{center}
\small
\begin{tabular}{rrrrr}
\hline
ambang & recall & FPR & presisi & akurasi \\
\hline
3,0000 & 0,450 & 0,00132 & 0,254 & 0,99813 \\
2,8000 & 0,510 & 0,00263 & 0,162 & 0,99688 \\
2,6000 & 0,550 & 0,00467 & 0,105 & 0,99488 \\
2,3263 & 0,620 & 0,01014 & 0,058 & 0,98949 \\
2,0000 & 0,740 & 0,02269 & 0,032 & 0,97707 \\
1,5000 & 0,860 & 0,06729 & 0,013 & 0,93264 \\
\hline
\end{tabular}
\end{center}

Pada kurva ROC, baris-baris itu bergerak dari $(0{,}001;0{,}45)$ ke $(0{,}067;0{,}86)$:
jejak yang naik tajam di dekat sumbu tegak, persis seperti yang diinginkan setiap
praktisi. AUC-nya $0{,}966448$. Tetapi kolom presisi mengatakan hal lain. Bahkan pada
ambang paling ketat, tempat hanya $132$ koneksi bersih ditandai, presisinya $0{,}254$:
tiga dari empat alarm tetap salah. Ketika recall dinaikkan ke $0{,}86$, presisinya
$0{,}013$.

Gambar~\ref{fig:presisi-recall-ambang} menggambar keenam baris tabel itu pada satu sumbu
ambang: recall menanjak sementara presisi menurun, dan tidak ada titik yang memberi keduanya
tinggi sekaligus.

\begin{figure}[htbp]
\centering
\begin{tikzpicture}
\begin{axis}[
  width=0.78\linewidth, height=6.2cm,
  x dir=reverse, xmin=1.45, xmax=3.05,
  ymin=0.35, ymax=0.95, ytick={0.4,0.5,0.6,0.7,0.8,0.9},
  xlabel={ambang $t$ (ke kanan: ambang dilonggarkan)}, ylabel={recall},
  axis y line*=left, axis x line*=bottom,
  tick label style={font=\footnotesize}, label style={font=\footnotesize},
  grid=major, grid style={layergrey!15},
]
\addplot[accent,thick,mark=*,mark size=1.8pt] coordinates {
  (3.0,0.450) (2.8,0.510) (2.6,0.550) (2.3263,0.620) (2.0,0.740) (1.5,0.860)};
\node[accent,font=\footnotesize,anchor=south west] at (axis cs:2.62,0.665) {recall};
\end{axis}
\begin{axis}[
  width=0.78\linewidth, height=6.2cm,
  x dir=reverse, xmin=1.45, xmax=3.05,
  ymin=0, ymax=0.30, ytick={0,0.1,0.2,0.3},
  axis y line*=right, axis x line*=none, ylabel={presisi (skala kanan)},
  tick label style={font=\footnotesize}, label style={font=\footnotesize},
]
\addplot[layerteal,thick,mark=square*,mark size=1.6pt] coordinates {
  (3.0,0.254) (2.8,0.162) (2.6,0.105) (2.3263,0.058) (2.0,0.032) (1.5,0.013)};
\node[layerteal,font=\footnotesize,anchor=north east] at (axis cs:2.78,0.232) {presisi};
\end{axis}
\end{tikzpicture}
\caption{Presisi dan recall pada enam ambang dari tabel sebelumnya, digambar pada satu
sumbu ambang yang diputar sehingga ke kanan ambang dilonggarkan. Recall menanjak dari
$0{,}450$ pada $t=3{,}0$ ke $0{,}860$ pada $t=1{,}5$, sedangkan presisi menurun dari
$0{,}254$ ke $0{,}013$: kedua kurva bergerak berlawanan arah dan tidak ada satu titik pun
yang memberi keduanya tinggi. Skala presisi ada di sisi kanan, $0$ sampai $0{,}30$; pada
$t=2{,}3263$, ambang yang dipakai simulasi, presisi $0{,}058$ dengan recall $0{,}620$.}
\label{fig:presisi-recall-ambang}
\end{figure}

Kedua bacaan itu tidak bertentangan: keduanya mengukur hal berbeda.

\begin{proposisi}
Kurva ROC dan AUC hanya bergantung pada sebaran skor bersyarat $S\mid y=1$ dan
$S\mid y=0$, sehingga tidak berubah bila kelaziman $\pi$ diubah. Presisi pada ambang
tetap berubah, dan mengikuti
\[
\operatorname{prec}(t)=\frac{\pi\,\mathrm{TPR}(t)}
{\pi\,\mathrm{TPR}(t)+(1-\pi)\,\mathrm{FPR}(t)} .
\]
\end{proposisi}

\begin{proof}[Bukti]
Menurut definisi, $\mathrm{TPR}$ dan $\mathrm{FPR}$ adalah peluang bersyarat pada $y$,
jadi keduanya tidak memuat $\pi$; kurva yang tersusun dari keduanya pun tidak memuat
$\pi$. Untuk presisi, hukum peluang total memberi
$P(\hat y=1)=\pi\,\mathrm{TPR}+(1-\pi)\mathrm{FPR}$, sedangkan
$P(y=1,\hat y=1)=\pi\,\mathrm{TPR}$. Aturan Bayes dari Bab 17 memberi bentuk di atas.
\end{proof}

Akibatnya dapat diukur. Bila kelaziman diubah menjadi setengah tanpa mengubah
sebaran skor bersyarat, AUC tinggal $0{,}967061$ (selisih $0{,}0006$ hanyalah galat
Monte Carlo), sedangkan \emph{average precision} melompat dari $0{,}287636$ menjadi
$0{,}967544$. Pada kelaziman satu banding seribu garis dasarnya $\pi=0{,}001$, jadi
$0{,}288$ masih jauh lebih baik daripada menebak, tetapi jauh lebih dekat kepada
kenyataan operasional daripada $0{,}966$.

Aturan praktisnya sederhana: bila kelas positif langka, laporkan kurva presisi--recall,
bukan hanya ROC, dan sebutkan kelazimannya. ROC tetap berguna untuk membandingkan dua model
pada populasi yang sama sebab ia tidak tercemar komposisi; ia menyesatkan begitu dipakai
menyimpulkan kegunaan pada populasi lain atau pada ambang tertentu. AUC juga buta terhadap
kalibrasi: model yang skornya dua kali terlalu besar memiliki AUC yang sama persis, sebab
urutannya tidak berubah. Dua cacat itu membawa kita ke bagian berikutnya.

\section{Kalibrasi: apakah angka itu berarti peluang}
\tingkat{menengah}

Operator yang menerima alarm perlu tahu satu hal yang tidak dijawab presisi maupun AUC:
bila model menulis $0{,}8$, apakah delapan dari sepuluh kasus seperti itu memang positif.
Kalimat sehari-harinya: kalibrasi memeriksa apakah ramalan cuaca yang mengatakan ``tujuh
puluh persen hujan'' memang benar pada tujuh dari sepuluh hari seperti itu.

\kamus{kalibrasi}{calibration}{kesesuaian antara angka yang dijanjikan model dan
kenyataan yang benar-benar terjadi pada kelompok dengan janji yang sama}

\begin{definisi}[Kalibrasi sempurna]
\index{kalibrasi}
Skor $S\in[0,1]$ disebut terkalibrasi sempurna bila
$P(y=1\mid S=s)=s$ untuk setiap $s$ yang mungkin. Pada data terhingga, kalibrasi diperiksa
dengan mengelompokkan skor ke dalam selang, lalu membandingkan rataan skor di selang itu
dengan laju positif yang teramati. Selisih mutlaknya, ditimbang ukuran selang, disebut
\emph{expected calibration error}.
\end{definisi}

Contoh kerja dengan $1000$ prediksi memperlihatkan bentuk pemeriksaannya.

\begin{center}
\small
\begin{tabular}{lrrrr}
\hline
selang skor & banyak & positif & rataan skor & laju positif \\
\hline
$[0{,}5;0{,}6)$ & 40 & 21 & 0,55 & 0,525 \\
$[0{,}6;0{,}7)$ & 60 & 36 & 0,65 & 0,600 \\
$[0{,}7;0{,}8)$ & 120 & 90 & 0,75 & 0,750 \\
$[0{,}8;0{,}9)$ & 200 & 176 & 0,85 & 0,880 \\
$[0{,}9;1{,}0]$ & 580 & 551 & 0,95 & 0,950 \\
\hline
\end{tabular}
\end{center}

Selisih terbesar ada di selang kedua, $0{,}05$, dan kesalahan kalibrasi keseluruhannya
\[
\tfrac{40}{1000}(0{,}025)+\tfrac{60}{1000}(0{,}050)+\tfrac{120}{1000}(0)
+\tfrac{200}{1000}(0{,}030)+\tfrac{580}{1000}(0)=0{,}010 .
\]
Arah kesalahannya berbeda antar selang, tetapi besarnya kecil: satu poin persentase. Pemeriksaan silangnya gratis: banyaknya
contoh di kelima selang berjumlah $1000$, dan positifnya $874$.

Kalibrasi diukur oleh galat kuadrat, dan di situ ada identitas yang rapi.

\begin{proposisi}[Brier untuk skor terkalibrasi]
\index{brier}
Misalkan $S$ terkalibrasi sempurna dan $B=\mathbb{E}[(S-y)^2]$ adalah skor Brier. Maka
\[
B=\mathbb{E}\bigl[S(1-S)\bigr].
\]
Khususnya skor konstan $S=\pi$ terkalibrasi sempurna dengan $B=\pi(1-\pi)$, dan skor
konstan itu adalah peramal terkalibrasi yang paling tidak tajam.
\end{proposisi}

\begin{proof}[Bukti]
Bersyarat pada $S=s$, peubah $y$ bernilai $1$ dengan peluang $s$. Maka
\[
\mathbb{E}\bigl[(s-y)^2\mid S=s\bigr]
=s(s-1)^2+(1-s)s^2=s(1-s)\bigl[(1-s)+s\bigr]=s(1-s).
\]
Mengambil harapan atas $S$ memberi bentuk pertama. Untuk $S=\pi$ konstan, syarat
kalibrasi menuntut $P(y=1)=\pi$, yang memang benar, dan rumusnya memberi $\pi(1-\pi)$.
\end{proof}

Identitas itu menjelaskan dua hal: kalibrasi saja tidak cukup, sebab skor konstan $\pi$
terkalibrasi sempurna tetapi tidak membedakan contoh; dan untuk kejadian langka garis
dasar Brier sangat rendah, sehingga model yang tampak baik pada akurasi dapat kalah telak
pada Brier. Pada simulasi kita $\pi(1-\pi)=0{,}000999$, sedangkan label keras model memberi
Brier sama dengan laju galatnya, $1051/100\,000=0{,}010510$. Skor keterampilan Brier,
$1-B/B_{\mathrm{garis\ dasar}}$, bernilai $-9{,}52$: model itu sepuluh kali lebih buruk
daripada meramalkan kelaziman konstan.

Angka negatif itu bukan keanehan: label keras membuang informasi ketidakpastian, dan pada
kejadian langka informasi itulah yang paling berharga. Model yang mengeluarkan $0{,}06$
sedang mengatakan sesuatu yang benar, dan laporan yang hanya memuat akurasi membuang
kalimat itu.

\section{Ketidakpastian metrik}
\tingkat{lanjut}

Semua angka di atas dihitung dari sejumlah contoh yang terbatas, jadi semuanya bergoyang
bila datanya diganti. Kalimat sehari-harinya: sebuah metrik tanpa selang adalah satu
tembakan tanpa keterangan seberapa lebar sasarannya.

Bab 20 sudah menurunkan interval Wilson dan membuktikan sifatnya, jadi kita tidak
mengulangnya di sini. Yang perlu diingat adalah bentuknya untuk proporsi $k/n$:
\[
\frac{\hat p+\dfrac{z^2}{2n}\pm
z\sqrt{\dfrac{\hat p(1-\hat p)}{n}+\dfrac{z^2}{4n^2}}}{1+\dfrac{z^2}{n}},
\qquad z=1{,}96 .
\]

\begin{untukanak}
Dua penjaga gawang berlatih menepis penalti. Yang pertama menepis sembilan dari sepuluh
tendangan. Yang kedua menepis sembilan ratus dari seribu tendangan. Koran menulis judul
yang sama untuk keduanya: sembilan puluh persen.

Tetapi kita merasa kedua kalimat itu tidak sama kuatnya, dan perasaan itu benar. Penjaga
pertama baru menghadapi sepuluh tendangan. Bisa jadi ia memang hebat, bisa jadi ia hanya
beruntung sepuluh kali berturut-turut. Penjaga kedua sudah menghadapi seribu tendangan,
dan keberuntungan sepanjang itu hampir mustahil. Angka yang sama, kekuatan yang berbeda.

Pertanyaan yang membuat penasaran: kalau kita hanya boleh menulis satu kalimat tentang
penjaga pertama, apakah ``sembilan puluh persen'' masih boleh ditulis? Boleh, asalkan
kalimatnya dilengkapi. Di bab ini kita akan tahu mengapa angka yang sama berarti ``antara
lima puluh sembilan dan sembilan puluh delapan persen'' pada sepuluh tendangan, tetapi
hanya ``antara delapan puluh delapan dan sembilan puluh dua persen'' pada seribu tendangan,
dan mengapa memperbanyak tendangan seratus kali lipat hanya mempersempit jawaban menjadi
sepersepuluh.
\end{untukanak}

Tabel berikut memakai rumus itu pada akurasi, dengan $z=1{,}96$.

\begin{center}
\small
\begin{tabular}{lrrr}
\hline
data & $\hat p$ & selang Wilson & lebar \\
\hline
$9/10$ & 0,900 & $(0{,}5959;\ 0{,}9821)$ & 0,3863 \\
$900/1000$ & 0,900 & $(0{,}8798;\ 0{,}9171)$ & 0,0372 \\
$10/10$ & 1,000 & $(0{,}7225;\ 1{,}0000)$ & 0,2775 \\
$0/1000$ & 0,000 & $(0{,}0000;\ 0{,}0038)$ & 0,0038 \\
\hline
\end{tabular}
\end{center}

Baris pertama dan kedua menjawab dua penjaga gawang tadi. Angka titiknya
sama, $0{,}9$, tetapi selang pertama selebar $0{,}3863$ dan yang kedua $0{,}0372$:
sepuluh kali lebih sempit, bukan seratus kali.

Baris ketiga adalah kasus yang paling sering muncul: tidak ada satu pun kesalahan
dalam $10$ percobaan, lalu ditulis ``akurasi $100$ persen''. Selang Wilson yang
jujur, $(0{,}7225;1)$, selebar $0{,}2775$: laju galatnya masih boleh $27{,}75$
persen. Di ujung lain, nol kejadian dalam $1000$ percobaan memberi batas atas
$0{,}003827$, sedangkan aturan tiga dari Bab 20 memberi $0{,}003$; keduanya boleh,
asalkan prosedurnya disebut.

Membandingkan dua model dengan memeriksa apakah selang masing-masing bertumpuk adalah
kebiasaan yang salah: keduanya diuji pada data yang sama, dan ketergantungan itu justru
menajamkan perbandingan.

\begin{proposisi}[Selisih dua model berpasangan]
Misalkan $b$ adalah banyak contoh yang benar oleh model A dan salah oleh model B, dan $c$
sebaliknya, pada $n$ contoh yang sama. Selisih akurasinya $d=(b-c)/n$ dengan variansi
\[
\operatorname{Var}(d)=\frac{1}{n}\Bigl(p_b+p_c-(p_b-p_c)^2\Bigr),
\]
sehingga galat baku tertaksirnya $\sqrt{b+c-(b-c)^2/n}\,/\,n$.
\end{proposisi}

\begin{proof}[Bukti]
Setiap contoh masuk salah satu dari empat sel: benar keduanya, salah keduanya, benar hanya
A, benar hanya B. Karena contohnya bebas, $(b,c)$ menyebar multinomial dengan peluang
$p_b,p_c$. Maka $\operatorname{Var}(b-c)=n\bigl(p_b+p_c-(p_b-p_c)^2\bigr)$, dan membagi
dengan $n^2$ memberi rumus di atas.
\end{proof}

Pada $b=37$, $c=15$, dan $n=1000$, selisihnya $0{,}022$ dengan galat baku $0{,}007177$,
sehingga selang $95$ persennya $(0{,}007932;0{,}036068)$, seluruhnya di atas nol, dan uji
McNemar eksak memberi nilai-p $0{,}003185$. Perlakuan naif yang menganggap kedua model
bebas memberi selang $(-0{,}004296;0{,}048296)$ yang memuat nol. Kedua model karena itu
dapat dinyatakan berbeda, asalkan perbandingannya berpasangan dan pengujiannya memakai
selisih, bukan dua selang terpisah.

Terakhir, ukuran uji menentukan selisih sekecil apa yang dapat dilihat. Untuk mendeteksi
selisih $\delta$ pada taraf $\alpha$ dan daya $1-\beta$,
\[
n\approx\frac{(z_{1-\alpha/2}+z_{1-\beta})^2\,\bigl(p_1(1-p_1)+p_0(1-p_0)\bigr)}
{\delta^2}.
\]
Pada $p_0=p_1=0{,}9$, taraf $0{,}05$, dan daya $0{,}8$, selisih satu poin menuntut
$n\approx14\,128$ contoh per model, selisih dua poin menuntut $3532$. Kumpulan uji seribu
contoh tidak dapat membedakan dua model yang terpaut satu poin.

\section{Kebocoran data dan validasi silang bersarang}
\tingkat{lanjut}

Semua perhitungan di atas mengandaikan satu hal yang tidak boleh dilanggar: angka uji
diukur pada data yang tidak pernah dilihat model, dan tidak pernah dilihat pula oleh
prosedur yang memilih model. Kalimat sehari-harinya: nilai ujian hanya berarti bila
soalnya belum pernah dibocorkan, termasuk oleh guru yang menyusun soalnya.

\kamus{kebocoran data}{data leakage}{jawaban ujian yang tanpa sengaja ikut masuk ke
bahan belajar, sehingga nilai menjadi tinggi tanpa kepandaian}

\begin{definisi}[Kebocoran data]
\index{kebocoran data}
Kebocoran terjadi bila informasi yang tidak akan tersedia pada saat model dipakai masuk
ke dalam pelatihan, pemilihan model, atau pelaporan, sehingga angka evaluasi mengukur
kemampuan mengingat, bukan kemampuan menggeneralisasi.
\end{definisi}

Bentuknya bermacam-macam, dan yang paling sering tidak terlihat. Menghitung rataan
atau median prapemrosesan dari seluruh data termasuk data uji adalah kebocoran:
statistik itu seharusnya dari latih saja. Memilih $50$ fitur terbaik dengan korelasi
terhadap label sebelum memisahkan latih dan uji juga kebocoran, sebab penyaring
fitur sudah melihat label uji. Menyetel hiperparameter pada kumpulan uji lalu
melaporkan angka terbaiknya adalah kebocoran yang paling sering dilakukan dengan
sadar. Pada data berkelompok, pemisahan acak menempatkan kelompok yang sama di kedua
sisi; pada deret waktu, ia membocorkan masa depan ke masa lalu. Yang paling sulit
dideteksi: fitur yang merupakan akibat dari label, misalnya kolom ``tindakan yang
diambil'' yang diisi setelah insiden diketahui.

Kebocoran pemilihan fitur dapat diukur, dan ukurannya mengejutkan. Ambil $200$ contoh
berlabel acak dan $2000$ fitur derau normal, tanpa sinyal apa pun. Pilih $20$ fitur dengan
korelasi mutlak terbesar terhadap label, lalu validasi silang lima lipatan dengan
pengklasifikasi centroid terdekat. Bila pemilihan
dilakukan sekali pada seluruh data sebelum lipatan dibentuk, akurasinya $0{,}7800$; bila
pemilihan diulang di dalam tiap lipatan memakai baris latihnya saja, $0{,}4600$. Pada tiga
benih lain angkanya $0{,}775$ dan $0{,}415$, $0{,}820$ dan $0{,}545$: yang bocor selalu
tinggi, yang jujur selalu di sekitar $0{,}5$. Angka $0{,}78$ itu seluruhnya ilusi, dan
tidak satu pun baris kodenya salah tulis.

Protokolnya satu kalimat: pisahkan lebih dulu, baru hitung apa pun. Bagi data
menjadi latih, validasi, dan uji; hitung statistik prapemrosesan dari latih saja;
setel hiperparameter pada validasi; sentuh uji sekali, di akhir, untuk satu angka.
Bila validasi silang dipakai, prapemrosesan berada di dalam lipatan. Untuk data
berkelompok, lipatannya mengikuti kelompok; untuk data waktu, mengikuti urutan
waktu.

Gambar~\ref{fig:pemisahan-tiga} menggambar urutan itu sebagai tiga kotak dan satu panah
yang dilarang.

\begin{figure}[htbp]
\centering
\begin{tikzpicture}[every node/.style={font=\footnotesize},
  blok/.style={draw=ink,thick,rounded corners=2pt,align=center,
               minimum width=3.5cm,minimum height=1.5cm,inner sep=3pt}]
  \node[blok,fill=layerteal!12]   (latih) at (0.0,0) {\textbf{latih} $60\%$\\ hitung statistik};
  \node[blok,fill=layerorange!12] (val)   at (4.8,0) {\textbf{validasi} $20\%$\\ pilih hiperparameter};
  \node[blok,fill=layerindigo!12] (uji)   at (9.6,0) {\textbf{uji} $20\%$\\ sentuh sekali};
  \draw[-{Latex},ink,thick] (latih.east) -- (val.west)
    node[midway,above,font=\scriptsize] {latih ulang};
  \draw[-{Latex},ink,thick] (val.east) -- (uji.west)
    node[midway,above,font=\scriptsize] {kunci sebelum dilihat};
  \draw[-{Latex},accent,thick,dashed] (uji.south) .. controls (9.6,-1.85) and (0.0,-1.85)
    .. (latih.south);
  \draw[accent,very thick] (4.62,-1.42) -- (4.98,-1.06);
  \draw[accent,very thick] (4.62,-1.06) -- (4.98,-1.42);
  \node[accent,anchor=west,font=\scriptsize] at (5.12,-1.24) {dilarang};
  \node[anchor=north,align=center,text width=11.6cm] at (4.8,-1.98)
    {statistik uji, penyaring fitur, dan penyetelan hiperparameter tidak boleh menyeberang
     balik ke kotak latih};
\end{tikzpicture}
\caption{Tiga kumpulan data yang terpisah, dengan proporsi $60/20/20$; pada $1000$ koneksi
itu berarti $600$, $200$, dan $200$. Statistik prapemrosesan dihitung dari kotak latih,
hiperparameter dipilih di kotak validasi, dan kotak uji disentuh sekali di akhir. Panah
putus-putus dari uji kembali ke latih adalah kebocoran: satu panah itu cukup untuk
menaikkan akurasi dari $0{,}4600$ menjadi $0{,}7800$ pada data yang seluruhnya derau.}
\label{fig:pemisahan-tiga}
\end{figure}

Yang tersisa adalah satu masalah halus. Bila validasi silang dipakai sekaligus untuk
memilih hiperparameter dan untuk melaporkan mutu, angka yang dilaporkan adalah angka
terbaik dari banyak percobaan pada data yang sama, jadi ia terlalu optimistis. Jalan
keluarnya adalah dua lapis lipatan.

\kamus{validasi silang bersarang}{nested cross-validation}{dua lapis lipatan, lapisan luar
untuk menilai dan lapisan dalam untuk memilih, seperti dua ruang ujian yang terpisah}

\begin{definisi}[Validasi silang bersarang]
\index{validasi silang}
\index{validasi silang bersarang}
Misalkan $\Lambda$ himpunan hiperparameter. Data dibagi menjadi $K$ lipatan luar. Untuk
setiap lipatan luar $k$: latih pada sisa data dengan setiap $\lambda\in\Lambda$ memakai
$J$ lipatan dalam, pilih $\lambda_k$ yang terbaik, latih ulang pada seluruh bagian latih
lipatan $k$ dengan $\lambda_k$, lalu ukur pada lipatan $k$. Angka yang dilaporkan adalah
rata-rata $K$ angka itu. Penilaian dan pemilihan karena itu memakai data yang berbeda.
\end{definisi}

Biayanya $K$ kali lebih mahal, dan itulah sebabnya protokol ini sering dilewati; tanpa
lapisan luar, tidak ada satu angka pun dalam laporan yang dapat ditafsirkan sebagai
taksiran mutu pada data baru.

\section{Mengapa satu kali uji tidak cukup}
\tingkat{menengah}

Kalimat ``kami menguji model pada $1000$ contoh dan memperoleh akurasi $0{,}900$''
terdengar tuntas, padahal ia memuat tiga ketidakpastian. Kalimat sehari-harinya:
satu kali menimbang tidak memberi tahu apakah timbangannya bergoyang, barangnya
berubah, atau kita memilih hasil yang paling menyenangkan.

Yang pertama statistis dan sudah dihitung: pada $n=1000$ selang $95$ persen selebar
$0{,}037$, dan memisahkan dua model yang terpaut satu poin menuntut $14\,128$ contoh.
Kesimpulan ``A lebih baik'' dari selisih sekecil itu adalah kesimpulan tentang undian.

Yang kedua selektif, dan namanya \emph{winner's curse}. Melatih $20$ varian yang mutunya
identik lalu melaporkan yang terbaik berarti melaporkan maksimum dari $20$ goyangan: pada
simulasi dengan akurasi benar $0{,}90$ dan $n=1000$, rata-rata tiap varian $0{,}899987$
sedangkan rata-rata maksimumnya $0{,}917337$, jadi $1{,}7$ poin dari angka itu
keberuntungan. Bila varian saling bergantung, goyangannya lebih kecil, tetapi arahnya
selalu ke atas.

Yang ketiga prosedural: kumpulan uji yang dipakai berulang berhenti menjadi kumpulan uji.
Bila tiap perbandingan punya peluang $5$ persen menipu, peluang sekurang-kurangnya satu
dari $20$ perbandingan menipu adalah $1-0{,}95^{20}=0{,}6415$, lebih besar daripada peluang
benar.

Obatnya sepadan: laporkan selang, pisahkan data pemilihan dari data pelaporan, dan kunci
kumpulan uji sebelum eksperimen dimulai. Bila uji terlanjur dipakai berulang, angka
berikutnya disebut angka eksplorasi.

\section{Daftar periksa pelaporan}
\tingkat{dasar}

Tujuh baris ini memisahkan laporan yang dapat diperiksa dari yang hanya meyakinkan.
Kalimat sehari-harinya: resi tidak membuat barang sampai, tetapi tanpa resi tak ada yang
dapat ditelusuri.

\begin{enumerate}
  \item Ukuran tiap kelas dan kelazimannya, bukan hanya $n$ total.
  \item Keempat bilangan matriks konfusi, lalu presisi dan recall pada ambang yang
        benar-benar dipakai, beserta nilai ambangnya.
  \item Selang kepercayaan tiap metrik utama, dengan metode yang disebut; untuk $n$ kecil
        pakai Wilson, bukan Wald.
  \item Kurva presisi--recall bila kelas positif langka, beserta average precision dan
        garis dasarnya $\pi$.
  \item Kalibrasi: batas selang skor, banyak contoh per selang, dan kesalahan kalibrasi.
  \item Protokol pemisahan: apa yang dihitung dari data mana, berapa kali data uji
        disentuh, dan berapa varian dicoba; angka terbaik tanpa penyebut itu angka
        pilihan, bukan angka mutu.
  \item Sebaran hasil antar benih atau antar lipatan, beserta kode dan benih acaknya.
\end{enumerate}

Yang tidak perlu ada: akurasi tanpa kelaziman, AUC tanpa kurva, dan kata ``siap
dipakai''.

\begin{resep}
Empat langkah melaporkan hasil evaluasi yang tidak menyesatkan:
\begin{enumerate}
  \item Pisahkan data lebih dulu: latih, validasi, uji, beserta proporsi dan kunci
        pemisahannya, ditulis sebelum satu hitungan pun dijalankan.
  \item Hitung prapemrosesan dan pilih fitur hanya dari latih, setel hiperparameter hanya
        pada validasi, lalu simpan uji untuk satu kali sentuh.
  \item Laporkan keempat sel matriks konfusi beserta $n$, $\pi$, dan ambang yang dipakai,
        lalu presisi, recall, dan selang kepercayaan tiap metrik.
  \item Sebutkan banyak varian yang dicoba dan sebaran hasil antar benih atau antar
        lipatan, supaya pembaca tahu berapa derajat kebebasan yang sudah dipakai.
\end{enumerate}
Peringatan urutan: langkah 1 tidak boleh menyusul langkah 2 atau 3. Statistik yang dihitung
sebelum data dipisahkan sudah memindahkan informasi uji ke latih, dan setelah itu tidak ada
angka yang dapat memperbaikinya; seluruh laporan hanya mengukur ingatan.
\end{resep}

\section{Contoh hitung}
\tingkat{menengah}

Tiga hitungan berikut dapat diulang di kertas.

Pertama, matriks konfusi dari catatan kerja penulis sendiri, pada $35\,069$ koneksi uji:
$\mathrm{TP}=23\,279$, $\mathrm{FP}=997$, $\mathrm{FN}=590$, $\mathrm{TN}=10\,203$.
\begin{align*}
\operatorname{prec}&=\frac{23\,279}{24\,276}=0{,}958931, &
\operatorname{rec}&=\frac{23\,279}{23\,869}=0{,}975282,\\
F_1&=\frac{46\,558}{48\,145}=0{,}967037, &
\mathrm{acc}&=\frac{33\,482}{35\,069}=0{,}954746,
\end{align*}
dan akurasi seimbangnya $0{,}943132$. Satu hal di catatan itu tidak konsisten: satu
kalimat menyebut $11\,200$ serangan dan $23\,869$ normal, sedangkan matriksnya menempatkan
$23\,869$ pada baris positif. Aritmetikanya benar, penamaan kelasnya tertukar, dan
kelaziman harus dibaca dari matriks, bukan dari kalimat.

Kedua, AUC dari data kecil. Positif berskor $0{,}92$, $0{,}81$, $0{,}55$ melawan negatif
berskor $0{,}70$, $0{,}41$, $0{,}35$, $0{,}22$ memberi $12$ pasangan; $0{,}92$ dan $0{,}81$
mengalahkan keempatnya, $0{,}55$ mengalahkan tiga, jadi $11/12=0{,}916667$. Bila satu skor
positif sama dengan satu skor negatif, pasangan itu dihitung setengah dan hasilnya
$0{,}875$; \texttt{sklearn} memberi angka yang sama bila $y$ dan skornya diumpankan
sebagai vektor.

Ketiga, efek kelaziman pada presisi ketika sensitivitas $0{,}90$ dan spesifisitas $0{,}95$
dipertahankan. Dengan $\operatorname{prec}=\pi(0{,}9)/\bigl(\pi(0{,}9)+(1-\pi)(0{,}05)\bigr)$:

\begin{center}
\small
\begin{tabular}{rrrr}
\hline
$\pi$ & $\operatorname{prec}$ & alarm benar & alarm palsu \\
\hline
0,001 & 0,017699 & 1 & 56 \\
0,010 & 0,153846 & 1 & 5,5 \\
0,100 & 0,666667 & 2 & 1 \\
0,500 & 0,947368 & 18 & 1 \\
\hline
\end{tabular}
\end{center}

Kolom terakhir adalah kebalikan presisi dikurangi satu, yaitu banyak alarm palsu untuk
setiap alarm benar. Pada kelaziman satu banding seribu, satu alarm benar datang bersama
$55{,}5$ alarm palsu; pada kelaziman setengah, bersama $0{,}056$. Model yang sama, mutu
yang sama, dan pengalaman operator yang sama sekali berbeda.

\kertas{Ambil matriks $1000$ prediksi tadi dan hitung ulang kelima metriknya tanpa melihat
teks. Lalu jawab: berapa $\mathrm{FP}$ yang harus dihapus, tanpa mengubah $\mathrm{TP}$
dan $\mathrm{FN}$, agar presisi menjadi $0{,}900$? Petunjuk: presisi $62/(62+x)=0{,}9$
memberi $x=6{,}89$, jadi $\mathrm{FP}$ harus turun dari $17$ menjadi paling banyak $6$.
Periksa bahwa akurasi sesudah perubahan itu $0{,}956$.}

\section{Lab komputasi: satu angka, empat cara membacanya}
\tingkat{lanjut}

Seluruh angka di bab ini dihitung oleh satu program pendek, yang juga memperlihatkan dua
hal yang tak dapat diperlihatkan rumus: seberapa besar kebocoran menaikkan angka, dan
seberapa besar pemenang dari banyak varian dilebihkan. Kodenya dijalankan apa adanya
dengan \texttt{numpy} dan \texttt{scipy}.

\begin{lab}
\begin{verbatim}
import numpy as np
from math import sqrt
from scipy.stats import norm

z = norm.ppf(0.975)


def wilson(k, n):
    p, z2 = k / n, z * z
    mid = (p + z2 / (2 * n)) / (1 + z2 / n)
    half = z * sqrt(p * (1 - p) / n + z2 / (4 * n * n)) / (1 + z2 / n)
    return mid - half, mid + half


def roc_pr(spos, sneg):
    s = np.concatenate([spos, sneg])
    y = np.concatenate([np.ones(len(spos)), np.zeros(len(sneg))])
    o = np.argsort(-s)
    s, y = s[o], y[o]
    tp, fp = np.cumsum(y), np.cumsum(1 - y)
    tpr, fpr = tp / len(spos), fp / len(sneg)
    prec = tp / (tp + fp)
    auc = np.sum(np.diff(np.concatenate([[0], fpr])) *
                 (tpr + np.concatenate([[0], tpr[:-1]])) / 2)
    ap = np.sum(np.diff(np.concatenate([[0], tpr])) * prec)
    return auc, ap


rng = np.random.default_rng(20261025)
Nneg, Npos = 99900, 100
sneg = rng.normal(0.0, 1.0, Nneg)
spos = rng.normal(2.6, 1.0, Npos)
t = 2.3263
tp = int((spos > t).sum()); fp = int((sneg > t).sum())
fn, tn = Npos - tp, Nneg - fp
n = Nneg + Npos
print("TP=%d FP=%d FN=%d TN=%d" % (tp, fp, fn, tn))
print("akurasi %.6f  garis dasar %.6f" % ((tp + tn) / n, Nneg / n))
print("presisi %.6f  recall %.6f" % (tp / (tp + fp), tp / Npos))
print("akurasi seimbang %.6f" % ((tp / Npos + tn / Nneg) / 2))
print("brier label keras %.6f  brier konstan %.6f"
      % ((fp + fn) / n, 0.001 * 0.999))
a1, p1 = roc_pr(spos, sneg)
print("pi=0,001: AUC=%.6f  AP=%.6f" % (a1, p1))
a2, p2 = roc_pr(rng.normal(2.6, 1.0, 50000), rng.normal(0.0, 1.0, 50000))
print("pi=0,500: AUC=%.6f  AP=%.6f" % (a2, p2))
for k, m in ((9, 10), (900, 1000), (10, 10), (0, 1000)):
    lo, hi = wilson(k, m)
    print("%4d/%-4d (%.6f, %.6f) lebar %.6f" % (k, m, lo, hi, hi - lo))
b, c = 37, 15
d = (b - c) / 1000
se = sqrt(b + c - (b - c) ** 2 / 1000) / 1000
print("selisih %.6f  se %.6f  IK (%.6f, %.6f)"
      % (d, se, d - z * se, d + z * se))
rt = np.random.default_rng(11)
a = rt.binomial(1000, 0.90, (20000, 20)) / 1000
print("20 varian: rata2 %.6f  rata2 maks %.6f"
      % (a.mean(), a.max(axis=1).mean()))

rl = np.random.default_rng(4)
nn, dd, kk = 200, 2000, 20
X = rl.normal(size=(nn, dd))
y = rl.permutation(np.repeat([0, 1], nn // 2))
corr = np.array([abs(np.corrcoef(X[:, j], y)[0, 1]) for j in range(dd)])
sel_all = np.argsort(-corr)[:kk]


def cv(pick):
    idx = rl.permutation(nn)
    ok = 0
    for f in np.array_split(idx, 5):
        tr = np.setdiff1d(idx, f)
        cols = pick(tr)
        m1 = X[np.ix_(tr[y[tr] == 1], cols)].mean(0)
        m0 = X[np.ix_(tr[y[tr] == 0], cols)].mean(0)
        d1 = ((X[np.ix_(f, cols)] - m1) ** 2).sum(1)
        d0 = ((X[np.ix_(f, cols)] - m0) ** 2).sum(1)
        ok += (np.where(d1 < d0, 1, 0) == y[f]).sum()
    return ok / nn


def top_in(tr):
    c = np.array([abs(np.corrcoef(X[tr, j], y[tr])[0, 1]) for j in range(dd)])
    return np.argsort(-c)[:kk]


print("pemilihan fitur: bocor %.4f  jujur %.4f"
      % (cv(lambda tr: sel_all), cv(top_in)))
\end{verbatim}
\end{lab}

Keluaran yang sebenarnya, apa adanya:

\begin{verbatim}
TP=62 FP=1013 FN=38 TN=98887
akurasi 0.989490  garis dasar 0.999000
presisi 0.057674  recall 0.620000
akurasi seimbang 0.804930
brier label keras 0.010510  brier konstan 0.000999
pi=0,001: AUC=0.966448  AP=0.287636
pi=0,500: AUC=0.967061  AP=0.967544
   9/10   (0.595850, 0.982124) lebar 0.386274
 900/1000 (0.879848, 0.917091) lebar 0.037243
  10/10   (0.722467, 1.000000) lebar 0.277533
   0/1000 (0.000000, 0.003827) lebar 0.003827
selisih 0.022000  se 0.007177  IK (0.007932, 0.036068)
20 varian: rata2 0.899987  rata2 maks 0.917330
pemilihan fitur: bocor 0.7800  jujur 0.4600
\end{verbatim}

Empat hal perlu dibaca dari keluaran itu. Pertama, dua angka pada baris kedua harus muncul
berdekatan: $0{,}989490$ dan $0{,}999000$. Arah selisihnya menentukan seluruh tafsiran:
model yang diuji berada di bawah garis dasar.

Kedua, AUC hampir tidak bergerak, $0{,}966448$ melawan $0{,}967061$, sedangkan AP bergerak
dari $0{,}287636$ ke $0{,}967544$, faktor $3{,}4$: bukti numerik bagi proposisi sebelumnya.

Ketiga, lebar selang menyusut sebanding $1/\sqrt n$, dari $0{,}386274$ pada $n=10$ menjadi
$0{,}037243$ pada $n=1000$, dan $0/1000$ memberi batas atas $0{,}003827$.

Keempat, dua baris terakhir memuat dua penyakit yang tak terlihat pada satu angka: selisih
berpasangan $0{,}022$ dengan selang $(0{,}007932;0{,}036068)$ tidak memuat nol, padahal
$(0{,}8798;0{,}9171)$ dan $(0{,}8563;0{,}8969)$ bertumpuk lebar; dan pemenang dari $20$
varian identik dilaporkan $0{,}917337$ padahal mutu sebenarnya $0{,}899987$.

Program tidak membuktikan satu pun proposisi di bab ini; ia memeriksa klaim yang sudah
dibuktikan. Simulasi memakai sebaran yang kita andaikan, dan validasi silang yang jujur
pada data derau memberi $0{,}4600$, bukan $0{,}5000$: angka tanpa sinyal pun bergoyang.

\begin{untukdosen}
Urutan mengajar yang saya sarankan: mulai dari laporan vendor pada pembuka bab, tulis
$99$ persen dan $99{,}9$ persen di papan, lalu minta kelas memutuskan mana yang lebih
baik. Hampir seluruh kelas memilih detektor, dan dari kegelisahan itu matriks konfusi
muncul sebagai jawaban, bukan definisi yang dihafal. Setelah presisi dan recall ada,
barulah $F_1$ dan akurasi seimbang masuk. Kurva ROC sebaiknya digambar tangan untuk enam
titik pada tabel ambang, sebab di situ mahasiswa melihat sendiri bahwa titik yang tampak
bagus pada ROC dapat berpresisi $0{,}013$.

Dua miskonsepsi khas. Pertama, AUC dianggap ukuran mutu yang bebas konteks sehingga ambang
dipilih sembarang; gejalanya tampak ketika mahasiswa tidak dapat menjelaskan presisi pada
ambang yang dipakai. Kedua, validasi silang dianggap selalu jujur, padahal pemilihan fitur
atau standardisasi di luar lipatan sudah cukup menaikkan angka $30$ poin pada data yang
seluruhnya derau.

Aktivitas papan tulis: bagi kelas menjadi dua kelompok, satu menghitung presisi untuk
$\pi=0{,}001$ dan satu untuk $\pi=0{,}5$ dengan sensitivitas dan spesifisitas yang sama,
lalu suruh keduanya mempertahankan hasilnya. Kelompok pertama menemukan $0{,}018$ dan harus
menjelaskan mengapa model yang sama terasa berbeda. Tanda kelas siap lanjut: mereka dapat
menyebut satu keadaan di mana akurasi $0{,}99$ lebih buruk daripada tidak memakai model.
\end{untukdosen}

\section{Ringkasan}
\tingkat{dasar}

\begin{ringkasan}
\begin{itemize}
  \item Empat bilangan matriks konfusi berperan tidak simetris: presisi menjawab
        pertanyaan operator, recall pertanyaan perancang kebijakan, dan keduanya
        bergantung pada ambang; $F_1$ adalah rata-rata harmoniknya dan mengabaikan
        $\mathrm{TN}$.
  \item Akurasi garis dasar adalah $1-\pi$: pada $\pi=0{,}001$ model dengan $0{,}989490$
        kalah dari program yang selalu menjawab negatif ($0{,}999000$), sedangkan akurasi
        seimbang $0{,}804930$ lebih jujur sebab tidak bergantung pada komposisi.
  \item AUC sama dengan peluang skor positif acak melebihi skor negatif acak, jadi hanya
        bergantung pada sebaran skor bersyarat; ia tidak berubah ketika kelaziman diubah,
        sedangkan presisi dan average precision berubah, dari $0{,}288$ ke $0{,}968$.
  \item Kalibrasi memeriksa apakah skor $0{,}8$ berarti delapan dari sepuluh. Untuk skor
        terkalibrasi, Brier sama dengan $\mathbb{E}[S(1-S)]$; skor konstan $\pi$
        terkalibrasi sempurna tetapi tidak tajam. Label keras di sini memberi
        $0{,}010510$ melawan $0{,}000999$.
  \item Lebar selang menyusut sebanding $1/\sqrt n$: $0{,}3863$ pada $n=10$ dan $0{,}0372$
        pada $n=1000$. Untuk membandingkan dua model, pakai selisih berpasangan: $b=37$,
        $c=15$ memberi $(0{,}007932;0{,}036068)$, sedangkan perlakuan tak berpasangan
        memberi selang yang memuat nol.
  \item Satu kali uji tidak cukup, dan kebocoran menaikkan angka secara palsu: $0{,}7800$
        bila fitur dipilih di luar lipatan, $0{,}4600$ bila di dalamnya; pemenang dari
        $20$ varian identik dilebihkan $1{,}7$ poin; dan $20$ perbandingan pada taraf $5$
        persen memberi peluang $0{,}6415$ bahwa salah satunya menipu.
\end{itemize}
\end{ringkasan}

\section{Latihan}
\tingkat{menengah}

\begin{latihan}
\begin{enumerate}
  \item Nilai klaim ini: ``Model A mencapai akurasi $99{,}1$ persen dan model B $98{,}7$
        persen pada kumpulan uji yang sama, jadi A lebih baik.'' Benar, salah, atau belum
        cukup bukti? Sebutkan bagian yang keliru dan dua angka
        tambahan yang memperkuat.
  \item Buktikan bahwa Brier untuk label keras $0/1$ sama dengan laju galat,
        $(\mathrm{FP}+\mathrm{FN})/n$. Mengapa kesamaan itu gagal untuk skor berupa
        peluang?
  \item Buktikan bahwa skor konstan $S=s$ meminimumkan Brier tepat ketika $s=\pi$, dengan
        nilai minimum $\pi(1-\pi)$.
  \item Diberikan $\mathrm{TP}=180$, $\mathrm{FP}=45$, $\mathrm{FN}=20$,
        $\mathrm{TN}=1755$. Hitung presisi, recall, $F_1$, akurasi, spesifisitas, dan
        akurasi seimbang sampai enam angka desimal, lalu periksa dengan
        \texttt{sklearn.metrics}.
  \item Hitung selang Wilson $95$ persen untuk $95/100$ dan $950/1000$. Bandingkan kedua
        lebar, lalu jelaskan mengapa perbandingan itu tidak boleh dipakai menyimpulkan
        model mana yang lebih baik.
  \item Hitung AUC dari lima positif berskor $0{,}95$, $0{,}88$, $0{,}70$, $0{,}60$,
        $0{,}40$ dan lima negatif berskor $0{,}75$, $0{,}65$, $0{,}50$, $0{,}30$,
        $0{,}10$. Tuliskan banyak pasangan benar, lalu periksa dengan
        \texttt{roc\_auc\_score}.
  \item Sebuah model memiliki sensitivitas $0{,}80$ dan spesifisitas $0{,}99$. Hitung
        presisinya pada $\pi=0{,}002$ dan $\pi=0{,}2$, lalu tentukan kelaziman minimum
        agar presisinya paling sedikit $0{,}5$.
  \item Anda memiliki $6000$ rekaman dari $300$ host, masing-masing sekitar $20$ rekaman,
        dan $40$ setelan hiperparameter. Rancang protokol evaluasi yang tidak bocor:
        sebutkan pembagian data, tempat prapemrosesan dihitung, tempat pemilihan
        hiperparameter, dan berapa kali uji disentuh.
  \item Nilai klaim ini: ``Karena AUC model ini $0{,}95$, model ini berguna pada ambang
        mana pun yang kami pilih.'' Benar, salah, atau belum cukup bukti? Sertakan satu
        perhitungan pendukung.
\end{enumerate}

Petunjuk untuk tiga butir tersulit. Butir 2: tulis $(s-y)^2$ untuk $s\in\{0,1\}$; nilainya
$1$ tepat pada kesalahan. Butir 3: turunkan $\mathbb{E}[(s-y)^2]=s^2-2s\pi+\pi$ terhadap
$s$, cari titik minimumnya, lalu pakai proposisi Brier di bab ini. Butir 7: presisinya
$0{,}8\pi/\bigl(0{,}8\pi+0{,}01(1-\pi)\bigr)$; selesaikan
$0{,}8\pi\ge0{,}5\bigl(0{,}8\pi+0{,}01(1-\pi)\bigr)$ terhadap $\pi$.
\end{latihan}

\begin{untukpeneliti}
Yang terbukti di bab ini: akurasi pengklasifikasi konstan adalah $1-\pi$; $F_1$ adalah
rata-rata harmonik presisi dan recall dengan $F_1\le(P+R)/2$; AUC sama dengan peluang
peringkat benar beserta bentuk Mann--Whitney-nya; ROC dan AUC invarian terhadap kelaziman
sedangkan presisi tidak; untuk skor terkalibrasi Brier sama dengan $\mathbb{E}[S(1-S)]$;
dan rumus variansi selisih dua proporsi berpasangan.

Yang populer tetapi bukan teorema: bahwa ``AUC $0{,}9$ berarti model bagus''; bahwa ambang
$0{,}5$ netral; bahwa validasi silang selalu jujur; dan bahwa selang yang bertumpuk berarti
tidak ada perbedaan. Tiga yang terakhir salah secara matematis, bukan hanya tidak lengkap.
\emph{Average precision} juga bukan taksiran tak bias bagi luas kurva presisi--recall
populasi: ia ringkasan dari kurva teramati, dan nilainya bergantung pada kelaziman sampel.

Sumber primer: Wasserman, \emph{All of Statistics}, Bab 6, Bab 9--10, Bab 13 bagian
Model Selection, dan Bab 22 bagian ``Assessing Error Rates and Choosing a Good Classifier''.
Goodfellow, Bengio, dan Courville, \emph{Deep Learning}, Bab 11 bagian ``Performance
Metrics'', memberi contoh penyakit satu per sejuta dan gagasan \emph{coverage}. Athalye,
Carlini, dan Wagner (2018) menunjukkan bahwa evaluasi yang lemah menghasilkan klaim
keamanan yang salah: tujuh dari sembilan pertahanan tanpa sertifikat pada ICLR 2018
bergantung pada gradien yang dikaburkan, dan enam tembus setelah evaluasinya diperbaiki.

Batas keberlakuan: seluruh simulasi memakai sebaran normal yang diandaikan, jadi ia
mengukur perilaku prosedur, bukan kebenaran model. Selang Wilson adalah aproksimasi
skor, dan McNemar mengandaikan pasangan tak sepakat simetris di bawah hipotesis nol.
Hasil kebocoran bergantung pada $n$, $d$, dan banyaknya fitur terpilih: $0{,}78$
bukan tetapan alam, melainkan satu titik pada kurva yang naik bersama $d/n$.

Pertanyaan terbuka: selang jujur untuk metrik dari lipatan yang saling bergantung;
pengujian kalibrasi tanpa asumsi selang skor seragam; dan pelaporan ketidakpastian ketika
kumpulan uji dipilih secara adaptif.
\end{untukpeneliti}

\peta{Bab 27 pada Wasserman, \emph{All of Statistics}: Bab 6 (penaksir, bias, galat baku),
Bab 9--10 (interval dan pengujian), Bab 13 bagian Model Selection, dan Bab 22 bagian
``Assessing Error Rates and Choosing a Good Classifier''. Interval Wilson dan aturan tiga
dirujuk dari Bab 20 buku ini serta dari Wilson (1927) dan Jovanovic--Levy (1997). Contoh
akurasi yang menipu mengikuti Goodfellow, Bengio, dan Courville, \emph{Deep Learning},
Bab 11; batas evaluasi pertahanan adversarial mengikuti Athalye, Carlini, dan Wagner (2018),
dibahas di Bab 27 dan 28. Aturan Bayes untuk presisi dibuktikan di Bab 16; pemisahan data
di Bab 24.}
